The rule for adding angular momenta obtained in \S~31 determines the possible values of the total angular momentum of a system composed of two particles (or more complex constituents) whose angular momenta are $j_1$ and $j_2$\footnote{Strictly, throughout what follows we shall tacitly mean a system composed of constituents whose mutual interaction is weak enough for their angular momenta to be treated as conserved in the first approximation. The results presented below apply, of course, not only when the total angular momenta of two particles (or systems) are added, but also when the orbital angular momentum and the spin of a single system are added, provided that the spin--orbit coupling is sufficiently weak.}. In fact, this rule is intimately connected with the way wave functions behave under rotations in space, and it follows directly from the properties of spinors.
